% Options for packages loaded elsewhere
\PassOptionsToPackage{unicode}{hyperref}
\PassOptionsToPackage{hyphens}{url}
\documentclass[
  ignorenonframetext,
]{beamer}
\newif\ifbibliography
\usepackage{pgfpages}
\setbeamertemplate{caption}[numbered]
\setbeamertemplate{caption label separator}{: }
\setbeamercolor{caption name}{fg=normal text.fg}
\beamertemplatenavigationsymbolsempty
% remove section numbering
\setbeamertemplate{part page}{
  \centering
  \begin{beamercolorbox}[sep=16pt,center]{part title}
    \usebeamerfont{part title}\insertpart\par
  \end{beamercolorbox}
}
\setbeamertemplate{section page}{
  \centering
  \begin{beamercolorbox}[sep=12pt,center]{section title}
    \usebeamerfont{section title}\insertsection\par
  \end{beamercolorbox}
}
\setbeamertemplate{subsection page}{
  \centering
  \begin{beamercolorbox}[sep=8pt,center]{subsection title}
    \usebeamerfont{subsection title}\insertsubsection\par
  \end{beamercolorbox}
}
% Prevent slide breaks in the middle of a paragraph
\widowpenalties 1 10000
\raggedbottom
\AtBeginPart{
  \frame{\partpage}
}
\AtBeginSection{
  \ifbibliography
  \else
    \frame{\sectionpage}
  \fi
}
\AtBeginSubsection{
  \frame{\subsectionpage}
}
\usepackage{iftex}
\ifPDFTeX
  \usepackage[T1]{fontenc}
  \usepackage[utf8]{inputenc}
  \usepackage{textcomp} % provide euro and other symbols
\else % if luatex or xetex
  \usepackage{unicode-math} % this also loads fontspec
  \defaultfontfeatures{Scale=MatchLowercase}
  \defaultfontfeatures[\rmfamily]{Ligatures=TeX,Scale=1}
\fi
\usepackage{lmodern}
\ifPDFTeX\else
  % xetex/luatex font selection
\fi
% Use upquote if available, for straight quotes in verbatim environments
\IfFileExists{upquote.sty}{\usepackage{upquote}}{}
\IfFileExists{microtype.sty}{% use microtype if available
  \usepackage[]{microtype}
  \UseMicrotypeSet[protrusion]{basicmath} % disable protrusion for tt fonts
}{}
\makeatletter
\@ifundefined{KOMAClassName}{% if non-KOMA class
  \IfFileExists{parskip.sty}{%
    \usepackage{parskip}
  }{% else
    \setlength{\parindent}{0pt}
    \setlength{\parskip}{6pt plus 2pt minus 1pt}}
}{% if KOMA class
  \KOMAoptions{parskip=half}}
\makeatother
\setlength{\emergencystretch}{3em} % prevent overfull lines
\providecommand{\tightlist}{%
  \setlength{\itemsep}{0pt}\setlength{\parskip}{0pt}}
% Slide layout defaults for warning-clean scientific decks.
\usepackage{etoolbox}
\IfFileExists{xurl.sty}{\usepackage{xurl}}{}
\IfFileExists{seqsplit.sty}{\usepackage{seqsplit}}{\newcommand{\seqsplit}[1]{#1}}
\protected\def\breakseq#1{\seqsplit{#1}}
\protected\def\breaktt#1{\begingroup\ttfamily\seqsplit{#1}\endgroup}
\setlength{\emergencystretch}{6em}
\tolerance=5000
\hbadness=10000
\hfuzz=1pt
\setlength{\tabcolsep}{2pt}
\AtBeginEnvironment{longtable}{\tiny\renewcommand{\arraystretch}{0.86}\setlength{\tabcolsep}{1pt}}
\AtBeginEnvironment{tabular}{\tiny\renewcommand{\arraystretch}{0.86}\setlength{\tabcolsep}{1pt}}

% Natbib and cross-reference fallbacks — slides don't load natbib
% or cleveref, but combined-PDF manuscript prose may emit these
% commands. The fallback renders citations as a bracketed key list
% and cross-references as detokenized labels so slides don't fail on
% undefined control sequences or raw underscores. \providecommand is
% a no-op if packages load later.
\providecommand{\citep}[1]{[#1]}
\providecommand{\citet}[1]{#1}
\providecommand{\citealp}[1]{#1}
\providecommand{\citeauthor}[1]{#1}
\providecommand{\citeyear}[1]{#1}
\providecommand{\cref}[1]{\texttt{\detokenize{#1}}}
\providecommand{\Cref}[1]{\texttt{\detokenize{#1}}}

\newtheorem{proposition}{Proposition}
\newtheorem{hypothesis}{Hypothesis}
\usepackage{bookmark}
\IfFileExists{xurl.sty}{\usepackage{xurl}}{} % add URL line breaks if available
\urlstyle{same}
% fallback for those not using the hyperref driver hyperxmp:
\makeatletter
\@ifundefined{xmpquote}{\newcommand{\xmpquote}[1]{#1}}{}
\makeatother
\hypersetup{
  hidelinks,
  pdfcreator={LaTeX via pandoc}}

\author{\texorpdfstring{}{}}
\date{}

\begin{document}

\begin{frame}[allowframebreaks]{Citation Network Topology
\label{sec:citation_network}}
\protect\phantomsection\label{citation-network-topology}
The intra-corpus citation network provides a structural view of how
Active Inference research is organized, identifying influential hub
papers, community structure, and patterns of citation isolation (Figure
\ref{fig:citation_network}).

\begin{figure}[htbp]
\centering
\includegraphics[width=0.9\textwidth]{../figures/citation_network.png}
\caption{Intra-corpus citation network (817 nodes, 2,176 edges). Node size reflects in-degree (number of intra-corpus citations received); highly cited foundational papers serve as nexus points connecting sub-domains.}
\label{fig:citation_network}
\end{figure}

\begin{block}{Network Density and Degree Distribution}
\protect\phantomsection\label{network-density-and-degree-distribution}
The intra-corpus citation network contains 817 nodes and 2,176 edges,
with a density of 0.33\% and 547 connected components. The average
in-degree of \(\approx 2.7\) indicates that most papers receive few
intra-corpus citations, consistent with the field's rapid expansion: the
majority of recent papers have not yet accumulated citations within the
corpus (Figure \ref{fig:degree_distribution}). Only 7.4\% of all
identified references (2,176 intra-corpus matches out of 29,323 total
reference entries) resolve to other papers within the corpus, reflecting
cross-source identifier mismatches and the field's engagement with a
broad external literature base. Community detection identifies clusters
via greedy modularity maximization \citep{clauset2004finding}.

\begin{figure}[htbp]
\centering
\includegraphics[width=0.7\textwidth]{../figures/degree_distribution.png}
\caption{In-degree distribution of the citation network. The power-law tail is characteristic of citation networks, with a small number of highly cited hubs.}
\label{fig:degree_distribution}
\end{figure}
\end{block}

\begin{block}{Connected Components and Citation Isolation}
\protect\phantomsection\label{connected-components-and-citation-isolation}
The high number of connected components (547 out of 817 nodes) reveals
that much of the corpus consists of citation-isolated papers---works
that neither cite nor are cited by other papers in the collection. A
single Giant Connected Component (GCC) typically dominates mature
scientific networks; here, with 547 components across 817 nodes, the GCC
contains a minority of nodes while the remainder form singletons or
small clusters of two to three papers. This is partially an artifact of
cross-source identifier mismatches, but it also reflects the field's
pattern of papers engaging with the FEP literature conceptually without
building explicit, graph-tractable citation chains. PageRank analysis
identifies highly influential papers, predominantly Friston's
foundational work \citep{friston2010free} and the AIF textbook
\citep{parr2022active}, which serve as nexus points linking otherwise
disconnected subgraphs.
\end{block}

\begin{block}{Network Summary}
\protect\phantomsection\label{network-summary}
\begin{table}[htbp]
\centering
\caption{Intra-corpus citation network summary statistics ($N = 817$ papers, 817 nodes). The low density and high component count reflect the field's rapid expansion and cross-source identifier mismatches.}
\label{tab:citation_network}
\begin{tabular}{ll}
\toprule
\textbf{Metric} & \textbf{Value} \\
\midrule
Nodes & 817 \\
Edges & 2,176 \\
Reference resolution rate & 7.4\% (2,176 / 29,323) \\
Connected components & 547 \\
Network density & 0.33\% \\
Mean in-degree & $\approx$ 2.7 \\
\bottomrule
\end{tabular}
\end{table}

The citation topology corroborates the field overview findings (RQ1,
RQ2): a small number of foundational papers---predominantly Friston's
free energy and active inference formulations---anchor a rapidly
expanding periphery of increasingly specialized work. The extremely low
density (0.33\%) corresponds to an epistemic stage of high
fragmentation, meaning that literature synthesis and cross-pollination
between specific sub-domains remain difficult. Theoretical influence
flows primarily through shared conceptual foundations (the hub nodes)
rather than through dense mutual citation across the periphery. As
metadata standardization improves and DOI adoption becomes universal
across preprint and journal ecosystems, re-running this pipeline should
yield substantially higher reference resolution rates and a more
connected graph, enabling finer-grained community detection and
tracking.
\end{block}
\end{frame}

\end{document}
